%----------------------------------------------------------------------------------------
%	PACKAGES AND OTHER DOCUMENT CONFIGURATIONS
%----------------------------------------------------------------------------------------

\documentclass[10pt,stdletter,dateno,sigleft]{newlfm} % Extra options: 'sigleft' for a left-aligned signature, 'stdletternofrom' to remove the from address, 'letterpaper' for US letter paper - consult the newlfm class manual for more options

%\usepackage{charter} % Use the Charter font for the document text
\usepackage{graphics}

\usepackage{standalone}

\usepackage{hyperref}
\usepackage[american]{babel}
\usepackage{csquotes}
\usepackage[backend=biber,style=authoryear,date=year,dashed=false,doi=false,isbn=false,url=false,arxiv=false]{biblatex}
\renewbibmacro{in:}{} % this should go immediately below calling biblatex 
%\DeclareLanguageMapping{american}{american-apa}
\addbibresource{/Users/hectorbahamonde/Bibliografia_PoliSci/library.bib} 
\hypersetup{
    bookmarks=true,         % show bookmarks bar?
    unicode=false,          % non-Latin characters in Acrobat's bookmarks
    pdftoolbar=true,        % show Acrobat's toolbar?
    pdfmenubar=true,        % show Acrobat's menu?
    pdffitwindow=true,     % window fit to page when opened
    pdfstartview={FitH},    % fits the width of the page to the window
    pdftitle={My title},    % title
    pdfauthor={Author},     % author
    pdfsubject={Subject},   % subject of the document
    pdfcreator={Creator},   % creator of the document
    pdfproducer={Producer}, % producer of the document
    pdfkeywords={keyword1} {key2} {key3}, % list of keywords
    pdfnewwindow=true,      % links in new window
    colorlinks=true,       % false: boxed links; true: colored links
    linkcolor=blue,          % color of internal links (change box color with linkbordercolor)
    citecolor=blue,        % color of links to bibliography
    filecolor=blue,      % color of file links
    urlcolor=blue           % color of external links
}

%% with this patch, the newlfm should work (BEGIN)
\usepackage{etoolbox}
\makeatletter
\patchcmd{\@zfancyhead}{\fancy@reset}{\f@nch@reset}{}{}
\patchcmd{\@set@em@up}{\f@ncyolh}{\f@nch@olh}{}{}
\patchcmd{\@set@em@up}{\f@ncyolh}{\f@nch@olh}{}{}
\patchcmd{\@set@em@up}{\f@ncyorh}{\f@nch@orh}{}{}
\makeatother
%% with this patch, the newlfm should work (END)


\newsavebox{\Luiuc}\sbox{\Luiuc}{\parbox[b]{1in}{\vspace{0.35in}
\includegraphics[width=1.2\linewidth]{/Users/hectorbahamonde/Job_Market/logo.png}}} % Company/institution logo at the top left of the page
\makeletterhead{Uiuc}{\Lheader{\usebox{\Luiuc}}}

\newlfmP{sigsize=50pt} % Slightly decrease the height of the signature field
\newlfmP{addrfromphone} % Print a phone number under the sender's address
\newlfmP{addrfromemail} % Print an email address under the sender's address
\PhrPhone{\texttt{p}} % Customize the "Telephone" text
\PhrEmail{\texttt{e}} % Customize the "E-mail" text

\lthUiuc % Print the company/institution logo

% email input begin
\newread\fid
\newcommand{\readfile}[1]% #1 = filename
{\bgroup
  \endlinechar=-1
  \openin\fid=#1
  \read\fid to\filetext
  \loop\ifx\empty\filetext\relax% skip over comments
    \read\fid to\filetext
  \repeat
  \closein\fid
  \global\let\filetext=\filetext
\egroup}
%\readfile{/Users/hectorbahamonde/Bibliografia_PoliSci/email.txt}
% email input end

\usepackage{color} % for colors
\usepackage{hyperref} % for links

\hypersetup{
  colorlinks = true,
  urlcolor = blue,
  pdfpagemode = UseNone
}

%----------------------------------------------------------------------------------------
%	YOUR NAME AND CONTACT INFORMATION
%----------------------------------------------------------------------------------------

\namefrom{\vspace{-3cm}Hector Bahamonde, PhD} % Name

\addrfrom{
{\vspace{-2cm}}\\ % Date
[12pt]
{\bf Hector Bahamonde, PhD}\\
{\color{blue}{\bf Senior Researcher}}\\
\emph{{\color{blue}INVEST Flagship Centre}}\\
{\color{blue}University of Turku}\\
Turku, Finland
}

\phonefrom{+358403621438} % Phone number

\emailfrom{\href{mailto:hibano@utu.fi}{hibano@utu.fi} \\ 
\texttt{w}: \href{http://www.hectorbahamonde.com}{www.HectorBahamonde.com}\\
\today}% Link to last version


%----------------------------------------------------------------------------------------
%	ADDRESSEE AND GREETING/CLOSING
%----------------------------------------------------------------------------------------

\greetto{\vspace{-1cm}Dear Rector of the University of Turku,} % Greeting text
%\closeline{\vspace{-1cm}\includegraphics[width=3cm]{/Users/hectorbahamonde/Administracion/signature.pdf} \hspace{-4cm}} % Closing text




%----------------------------------------------------------------------------------------

\begin{document}
\begin{newlfm}

%----------------------------------------------------------------------------------------
%	LETTER CONTENT
%----------------------------------------------------------------------------------------
 %introduction
\vspace{-0.3cm}I am writing to apply for the title of Docent at the Faculty of Social Sciences of the University of Turku. Currently, I am a \href{https://www.utu.fi/en/people/hector-bahamonde-norambuena}{Senior Researcher} at the INVEST Flagship Centre at the University of Turku. I also serve as the Research Director of the \href{https://www.utu.fi/en/investhub}{INVEST-Hub}, an interdisciplinary center about experimental and quantitative social research. I obtained my B.A. in Political Science at Catholic University of Chile. After receiving my PhD in Political Science from Rutgers University-New Brunswick, NJ, U.S. in May 2017, I spent one year as a post-doctoral fellow at Tulane University-New Orleans, LA, U.S. Before coming to the University of Turku, I was an Assistant Professor (tenure-track) at O$'$Higgins University, in Chile for two years. My long-term plan is to settle in Turku, and to build an academic career at the University of Turku, hopefully securing a tenure-track position in my field. Thus, the title of Docent, to me, is a stepping stone in my long-term development plan. In this statement I will summarize my accomplishments and future research and teaching plans.

{\bf Research and scientific activities}. Within Political Science, my primary subfield is the political economy of inequality, democratic development and clientelism. In simple, I am interested in how economic conditions---for example, inequality, economic growth, poverty, economic losses and gains, among others---affect democracy itself and people's attitudes toward democracy. Research methods are an important part of my research. Every piece I have published pays considerable attention to both the methodological and substantive aspects of the question at hand. Importantly, since I defended my dissertation in 2017, and in my humble opinion, I have been able to produce more than the equivalent amount of work corresponding to a dissertation. Next paragraph briefly describes (a selection of) my peer-reviewed publications.


In \href{https://doi.org/10.1177/1866802X1801000202}{\textcite{Bahamonde2018}}, and in the context of clientelism in Brazil, I asked whether political parties target individuals or groups. Using causal inference tools (i.e., matching), the paper finds that parties target individuals when they are identifiable and groups when parties need to rely on the spillover effects of clientelism. Next, in \href{https://doi.org/10.1016/j.ejpoleco.2021.102048}{\textcite{Bahamonde2021a}} we asked about the relationship between democracy, state capacity and economic inequality. Using time-series and cross-sectional methods, we find that higher levels of democracy and state capacity increase inequality overtime. We exploited a large dataset comprising all countries since the 1970s. After that, in \href{https://link.springer.com/article/10.1057/s41269-020-00174-4}{\textcite{Bahamonde2020a}}, I published the results of a list experiment fielded in the United States. Unlike other research questions that ask factual questions, in this piece I asked whether voters \emph{would} vote for a candidate in exchange for money. Thought clientelism is low in the U.S., that does not necessarily mean that voters \emph{would not} sell their vote to a political party. I find that 25\% of voters (it was a representative sample) would vote for a candidate in exchange for money. The same year we published in \href{https://doi.org/10.1016/j.electstud.2022.102497}{\textcite{Bahamonde2022b}} the results of a laboratory experiment. In particular, we asked when parties buy more votes, whether it is when they are winning or losing the election. We discovered that parties engage in clientelism more when they are winning the election because parties are more concerned with preventing a decline than increasing gains. Finally, in the context of a natural experiment in Finland, \textcite{Bahamonde:2023},\footnote{This paper is Revise and Resubmit. It is in its second round of revisions, and after seeing the reviews, we are quite confident that that paper will be published.} we discovered that candidates that look working class have less chances of winning an election, particularly female candidates.

As it can be seen, I have been quite active publishing in some of the top journals in the discipline. I'd take the liberty to argue as well that in the meantime, I had three children and also the COVID pandemic hit. My oldest son was just a few months old when he sat in my dissertation defense. Then, as it can be seen in my CV, several movings happened (within the U.S. and abroad as well). Then the pandemic hit, slowing me down even more, particularly with the collection of the the lab experiment paper. However, I'm a very determined person and managed to accomplish this list of publications. 



% ending
{More information, \href{http://www.hectorbahamonde.com/teaching/}{\texttt{syllabi}}, my \href{http://github.com/hbahamonde/Job_Market/raw/master/Bahamonde_Research_Statement.pdf}{\texttt{research}}, \href{http://github.com/hbahamonde/Job_Market/raw/master/Bahamonde_Teaching_Statement.pdf}{\texttt{teaching}} and \href{http://github.com/hbahamonde/Job_Market/raw/master/Bahamonde_Diversity_Statement.pdf}{\texttt{diversity}} statements, as well as other \href{http://www.hectorbahamonde.com/research/}{\texttt{papers}} are available on my website: \href{http://www.hectorbahamonde.com}{\texttt{www.HectorBahamonde.com}}. Thank you for considering my application. I look forward to hearing from you.\unskip}

% References
\newpage
\printbibliography


%----------------------------------------------------------------------------------------

\end{newlfm}
\end{document}

The academic portfolio is structured as follows:
#. Basic information
#- Personal and contact information
#- Education and degrees
#- Other relevant background information
#2. Research and academic activities
#- Experience in research and other scientific activities
- Research philosophy and ethics, primary fields of research
#- Most notable publications
- Research evaluations and awards
- Activity in the academic community
- Vision and personal development
- Other scientific qualifications
3. Teaching and supervision
- Experience of undergraduate and doctoral teaching and supervision
- Pedagogical approach and training
- Development of teaching skills and awards for teaching
- Participation in the assessment and development of teaching
- Strengths, development challenges and visions for own teaching
- Other qualifications related to teaching
4. Administrative tasks and other activities
- Administrative and management duties
- Positions within the field outside academia
- Societal activities and positions of trust
- Publications, presentations and the dissemination of scientific knowledge
through other means
- Other significant qualifications and tasks
